%% =================================================================
%% Chen et al. Measurement 264 (2026) 120300.
%% Reconstruction of page 8 (printed p. 8).
%% Contents: Fig. 10 (topographic-lines resolution study), Fig. 11
%% (gauge-block accuracy study), prose: tail of Section 4.1 (FEA
%% validity) and start of Section 4.2 (Scanning resolution).
%% =================================================================

\documentclass[11pt]{article}

\usepackage[margin=0.85in]{geometry}
\usepackage{amsmath, amssymb}
\usepackage{mathrsfs}
\usepackage{xcolor}
\usepackage{fancyhdr}

\pagestyle{fancy}
\fancyhf{}
\lhead{\small Z.~Chen et al.}
\rhead{\small Measurement 264 (2026) 120300}
\cfoot{\small 8 $\to$ \arabic{page}}
\renewcommand{\headrulewidth}{0.4pt}

\setcounter{page}{1}
\setcounter{equation}{1}
\setlength{\parskip}{6pt}
\setlength{\parindent}{0pt}

\begin{document}

% ---------- Fig. 10 placeholder ----------
\begin{center}
\fbox{\begin{minipage}[c][5.5cm][c]{0.95\textwidth}
\centering
\textbf{Fig. 10.}\ Scanning resolution of the tactile tomography
system.\\[4pt]
\textcolor{gray}{\footnotesize
\textbf{(a)} A model featuring topographic lines with line
spacings of 2, 1.5, 1, 0.7, 0.5, and 0.3 mm (all lines 2 mm wide,
8 mm long), 3D printed.\\[3pt]
\textbf{(b)} Topographic-lines model filled with a soft surface
layer of AB glue (photograph).\\[3pt]
\textbf{(c)} 3D imaging of compression-depth distribution when the
tactile tomography system scanned (b). Lines above 0.5 mm spacing
remain cleanly distinguishable in the reconstructed depth map;
lines at 0.3 mm spacing blur together. X 0--70 mm, Y 0--20 mm,
Z color scale 0.1--5.2 mm.]}
\end{minipage}}
\end{center}

\vspace{0.4em}

\textbf{Fig. 10.}\ Scanning resolution of the tactile tomography
system. (a) A model featuring topographic lines. (b) Topographic
lines model filled with a soft surface layer of AB glue. (c) 3D
imaging of compression depth distribution when the tactile
tomography system scanned the topographic lines model with a soft
surface layer.

\vspace{0.8em}

% ---------- Fig. 11 placeholder ----------
\begin{center}
\fbox{\begin{minipage}[c][6cm][c]{0.95\textwidth}
\centering
\textbf{Fig. 11.}\ Scanning accuracy of the tactile tomography
system.\\[4pt]
\textcolor{gray}{\footnotesize
\textbf{(a)} A gauge block with a length of 30 mm, a width of 9 mm,
and a height of 3 mm.\\[3pt]
\textbf{(b)} Gauge block filled with a soft surface layer of AB
glue; a 20 mm $\times$ 20 mm scanning area is marked on the top
surface.\\[3pt]
\textbf{(c)} 3D imaging of compression-depth distribution when the
tactile tomography system scanned (b). The gauge block reads as a
flat plateau at $Z \approx 0.1$ mm with a cliff down to
$Z \approx 2.4$ mm at the edges.\\[3pt]
\textbf{(d)} Scanning-accuracy distribution corresponding to the
compression-depth scan, color-coded 98.95\% (lightest) to 100.00\%
(darkest) across the 20$\times$20 mm scan area.]}
\end{minipage}}
\end{center}

\vspace{0.4em}

\textbf{Fig. 11.}\ Scanning accuracy of the tactile tomography
system. (a) A gauge block with a length of 30 mm, a width of 9 mm,
and a height of 3 mm. (b) A gauge block filled with a soft surface
layer of AB glue. (c) 3D imaging of compression depth distribution
when the tactile tomography system scanned the gauge block with a
soft surface layer. (d) Scanning accuracy distribution
corresponding to compression depth when the tactile tomography
system scanned the gauge block with a soft surface layer.

\vspace{0.6em}

% ---------- Prose ----------
compression and deformation to a certain displacement as the probe
tip presses down (Movie~S2). When the sample was scanned by the
probe tip with 676 scanning points under a constant pressure of 4
MPa, as shown in Fig.~9d and Movie~S3, the ``$+$'' shaped region
exhibit distinct stress responses compared to the surrounding
matrix, and the displacement value of these regions are less than
that of the surrounding matrix, reflecting mechanical interaction
differences induced by the internal structure. Fig.~9e presents
the simulated displacement field of the sample under probe loading.
According to the color-code displacement vectors, the ``$+$''
shaped region has lower displacement magnitudes than the surrounding
matrix, which aligns with the tactile tomography principle that
internal structures alter the effective thickness of the soft
surface layer, resulting in different distribution of the tactile
softness across the samples. In addition, this simulated result is
consistent with the experimental result (as shown in Fig.~9f),
confirming that the tactile tomography can effectively translate
mechanical interaction differences into detectable internal
structure maps.

\subsection*{4.2.\ Scanning resolution of the tactile tomography system}

Scanning resolution is the key performance parameter in determining
the detailed quality of scanned imaging. In this article, the x
stage, y stage, and z stage of the tactile tomography system can
achieve a minimum step displacement of 10, 10, and 0.1 $\mu$m,
respectively, which ensures the system has a high scanning
resolution. To investigate the scanning resolution of the tactile
tomography system, a model featuring topographic lines was
fabricated by 3D printing. As shown in Fig.~10a, these lines have
a length of 8 mm and a width of 2 mm in the x and y directions,
and are interspaced by 2.0, 1.5, 1.0, 0.7, 0.5, and 0.3 mm,
respectively. The model was then filled with a soft surface layer
of AB glue (Fig.~10b), and scanned by the tactile tomography system
under a constant pressure of 4 MPa. The scanning result is shown
in Fig.~10c, where the distances between the topographic lines in
the x and y directions above 0.5 mm can be distinguished, meaning
that the tactile tomography system possesses a minimum scanning
resolution of 0.5 mm in both the x and y directions. This scanning
resolution is typically limited by the diameter of the tip, which
is also 0.5 mm. Therefore, the scanning resolution in the x and y
directions can be improved by reducing the diameter of the tip.
But there is a trade-off between the scanning resolution and the
diameter of the tip, because a tip with a very small diameter will
pierce the soft surface layer easily. The scanning resolution in
the z direction is unaffected by the diameter of the tip, and
mainly depends on the minimum step displacement. Thus, the tactile
tomography system can achieve a high scanning resolution of 0.1
$\mu$m in the z direction, allowing the tactile tomography system
to recognize an object in more detail.

\subsection*{4.3.\ Scanning accuracy of the tactile tomography system}

Scanning accuracy is another key performance parameter for the
tactile tomography system, referring to its reliability and
repeatability.

\end{document}
